\def\thelstlisting{}  % potrzebne do hacka zmiany formatu numerowania listingów
\documentclass[a4paper,12pt]{report}
\usepackage{polski}
\usepackage[polish]{babel}
\usepackage[utf8x]{inputenc}
\usepackage{graphicx}
\usepackage{listings}
\usepackage{color}
\usepackage{csvsimple}
\usepackage{pgfplots}
\usepackage{lipsum}
\usepackage{titlesec, blindtext, color}
\usepackage{url}
\usepackage{indentfirst}
\graphicspath{{images/}}

% Polskie numerowanie:
\makeatletter
 \def\thechapter{\arabic{chapter}.}
 \def\thesection{\thechapter\arabic{section}.}
 %\def\thesection{\arabic{section}.}
 \def\thesubsection{\thesection\arabic{subsection}.}
 \def\thesubsubsection{\thesubsection\arabic{subsubsection}.}
\makeatother
\renewcommand{\thefigure}{\thechapter\arabic{figure}}
\renewcommand{\thetable}{\thechapter\arabic{table}}
% to jest hack, bo sie nie da zrobic normalnie :(
%\renewcommand{\thelstlisting}{\thechapter\arabic{lstlisting}.}
\makeatletter
\AtBeginDocument{%
  \renewcommand \thelstlisting
       {\ifnum \c@chapter>\z@ \thechapter\fi \@arabic\c@lstlisting}%
\renewcommand{\lstlistingname}{Listing}
\renewcommand{\tablename}{Tabela} 
\renewcommand{\figurename}{Rysunek} 
%\renewcommand{\lstlistlistingname}{} 
%\renewcommand{\listfigurename}{}
%\renewcommand{\listtablename}{}

}
\makeatother


%\newcommand{\sectionbreak}{\newpage}
%\let\stdsection\section
%\renewcommand\section{\newpage\stdsection}

\titleformat{\chapter}
{\normalfont\huge\bfseries}{\thechapter}{20pt}{\Huge}
\titleformat{\section}
{\normalfont\Large\bfseries}{\thesection}{1em}{}
\titleformat{\subsection}
{\normalfont\large\bfseries}{\thesubsection}{1em}{}
\titleformat{\subsubsection}
{\normalfont\normalsize\bfseries}{\thesubsubsection}{1em}{}
\titleformat{\paragraph}[runin]
{\normalfont\normalsize\bfseries}{\theparagraph}{1em}{}
\titleformat{\subparagraph}[runin]
{\normalfont\normalsize\bfseries}{\thesubparagraph}{1em}{}
\titlespacing*{\chapter} {0pt}{24pt}{24pt}
\titlespacing*{\section} {0pt}{12pt}{12pt}
\titlespacing*{\subsection} {0pt}{12pt}{12pt}
\titlespacing*{\subsubsection}{0pt}{12pt}{12pt}
\titlespacing*{\paragraph} {0pt}{3.25ex plus 1ex minus .2ex}{1em}
\titlespacing*{\subparagraph} {\parindent}{3.25ex plus 1ex minus .2ex}{1em}

\providecommand{\e}[1]{\ensuremath{\times 10^{#1}}}

\definecolor{mygreen}{rgb}{0,0.6,0}
\definecolor{mygray}{rgb}{0.5,0.5,0.5}
\definecolor{mymauve}{rgb}{0.58,0,0.82}

\lstset{
  backgroundcolor=\color{white},   % choose the background color
  basicstyle=\footnotesize,        % the size of the fonts that are used for the code
  breakatwhitespace=false,         % sets if automatic breaks should only happen at whitespace
  breaklines=true,                 % sets automatic line breaking
  captionpos=t,                    % sets the caption-position to bottom
  commentstyle=\color{mygreen},    % comment style
  escapeinside={\%*}{*)},          % if you want to add LaTeX within your code
  frame=tb,                    % adds a frame around the code
  keepspaces=true,                 % keeps spaces in text, useful for keeping indentation of code
  keywordstyle=\color{blue},       % keyword style
  language=c,                 % the language of the code
  numbers=left,                    % where to put the line-numbers; possible values are (none, left, right)
  numbersep=5pt,                   % how far the line-numbers are from the code
  numberstyle=\footnotesize\color{mygray}, % the style that is used for the line-numbers
  rulecolor=\color{black},         % if not set, the frame-color may be changed on line-breaks within not-black text
  showspaces=false,                % show spaces everywhere adding particular underscores; it overrides 'showstringspaces'
  showstringspaces=false,          % underline spaces within strings only
  showtabs=false,                  % show tabs within strings adding particular underscores
  stepnumber=1,                    % the step between two line-numbers. If it's 1, each line will be numbered
  stringstyle=\color{mymauve},     % string literal style
  tabsize=2,                       % sets default tabsize to 2 spaces
  caption=\lstname,                   % show the filename of files included with \lstinputlisting; also try caption instead of title
    inputencoding=utf8x, 
    extendedchars=\true,
    lineskip=1pt,
    literate={ą}{{\k{a}}}1
             {Ą}{{\k{A}}}1
             {ę}{{\k{e}}}1
             {Ę}{{\k{E}}}1
             {ó}{{\'o}}1
             {Ó}{{\'O}}1
             {ś}{{\'s}}1
             {Ś}{{\'S}}1
             {ł}{{\l{}}}1
             {Ł}{{\L{}}}1
             {ż}{{\.z}}1
             {Ż}{{\.Z}}1
             {ź}{{\'z}}1
             {Ź}{{\'Z}}1
             {ć}{{\'c}}1
             {Ć}{{\'C}}1
             {ń}{{\'n}}1
             {Ń}{{\'N}}1
}



\setlength{\paperheight}{297mm}
\setlength{\paperwidth}{210mm} 

\setlength{\hoffset}{0.00cm}
\setlength{\voffset}{0.00cm}

\setlength{\oddsidemargin}{1.00cm}
\setlength{\topmargin}{1mm}
\setlength{\textheight}{\paperheight}
\textheight 23 cm	%26
\voffset -1.5 cm		%-3
\setlength{\textwidth}{16cm}  
\hoffset -0.5 cm
\setlength{\marginparsep}{1mm}  
\setlength{\marginparwidth}{2cm}
\setlength{\footskip}{2.36cm}   
\setlength{\parindent}{10mm}

\linespread{1.5}
\setlength{\parindent}{10mm}

\hyphenpenalty=2000			% nie dziel wyrazów zbyt często
\clubpenalty=10000			% kara za sierotki
\widowpenalty=10000			% nie pozostawiaj wdów
\brokenpenalty=10000		% nie dziel wyrazów między stronami
% zakomentowane bo psuje listingi
%\exhyphenpenalty=999999		% nie dziel słów z myślnikiem 

% Alter some LaTeX defaults for better treatment of figures:
% See p.105 of "TeX Unbound" for suggested values.
% See pp. 199-200 of Lamport's "LaTeX" book for details.
%   General parameters, for ALL pages:
\renewcommand{\topfraction}{0.9}	% max fraction of floats at top
\renewcommand{\bottomfraction}{0.8}	% max fraction of floats at bottom
%   Parameters for TEXT pages (not float pages):
\setcounter{topnumber}{2}
\setcounter{bottomnumber}{2}
\setcounter{totalnumber}{4}     % 2 may work better
\setcounter{dbltopnumber}{2}    % for 2-column pages
\renewcommand{\dbltopfraction}{0.9}	% fit big float above 2-col. text
\renewcommand{\textfraction}{0.07}	% allow minimal text w. figs
%   Parameters for FLOAT pages (not text pages):
\renewcommand{\floatpagefraction}{0.7}	% require fuller float pages
% N.B.: floatpagefraction MUST be less than topfraction !!
\renewcommand{\dblfloatpagefraction}{0.7}	% require fuller float pages

\title{Analiza i porównanie popularncych, otwartoźródłowych bibliotek (framework'ów) do tworzenia aplikacji internetowych na podstawie prostej aplikacji.}

\author{Rafał Selewońko}
\date{\today}


\def\titlelinea{Wydział Informatyki}
\def\titlelineb{Katedra Oprogramowania}
\def\titlelinec{Obliczenia równoległe}
\def\titlelined{Temat: Program typu ping-pong}
\def\wykonawca{inż. Rafał Selewońko}
\def\promotor{mgr inż. Paweł Zabielski}
\def\promotorstring{Prowadzący: } 
\def\data{6 grudnia 2013}

\begin{document}

\begin{titlepage}
\begin{center}

% Upper part of the page. The '~' is needed because \\
% only works if a paragraph has started.
\includegraphics[width=0.5\textwidth]{./logo_PB}~\\[0.5cm]
\textsc{\LARGE \titlelinea}\\[0.5cm]
\textsc{\Large \titlelineb}\\[1.5cm]

% Title
\textsc{{ \LARGE  \titlelinec \\[0.5cm] }}
{ \large \bfseries \titlelined \\[1.5cm] }


% Author and supervisor
\begin{flushright} \large
Wykonawca: \textsc{\wykonawca}\\[1.5cm]
\end{flushright}


\begin{flushleft} \large
\promotorstring\textsc{\promotor}
\end{flushleft}

% Bottom of the page
\vfill
{\large \data}

\end{center}

%\newpage
%\thispagestyle{empty}
%\pagestyle{empty}
%\vspace*{15\baselineskip}
%\begin{center}
%\textit{42}
%\end{center}

\end{titlepage}

%\tableofcontents

\chapter{Zadanie}

Aplikacja MPI składa się z dwóch procesów: P0 oraz P1. Obydwa procesy 
wymieniają się komunikatami stosując wzorzec wymiany jak w grze w ping-ponga. P0 wysyła 
komunikat do P1 następnie P1 odsyła komunikat do P0. Ta sekwencja wymiany jest powtarzana 
odpowiednią (bardzo dużą) liczbę iteracji, tak aby program wykonywał się przez kilka sekund. Czas 
pracy programu jest obliczany przy pomocy funkcji MPI\_Walltime. Dzieląc ten czas przez liczbę 
iteracji pomnożoną przez dwa (ponieważ komunikat podróżuje w obie strony) otrzymujemy czas 
transmisji komunikatu w jedną stronę. Program powinien wykonywać szereg takich pomiarów dla 
rosnącej długości komunikatów, tak abyś mógł w stanie sporządzić wykres 
czas komunikatu=f(długość). Z kolei dzieląc długość komunikatu przez czas jego transmisji otrzymujemy przepustowość w bajtach/s 
dla danej długości komunikatu. Należy sporządzić również wykres przepustowość=f(długość). 
Program powinien implementować test ping-pong w dwóch wersjach: przy pomocy blokujących 
operacji MPI\_Send/MPI\_Recv oraz przy pomocy nieblokujących MPI\_Isend/MPI\_Recv wraz z 
MPI\_Wait. Wykonaj pomiar na klastrze i sporządź raport z wykonania zadania zawierający wydruk 
programu + obydwa wykresy. Porównaj otrzymane wyniki z wynikiem polecenia ping oraz z 
parametrami sieci Gigabit Ethernet. Porównaj wynik na czas transmisji komunikatu z modelem 
przedstawionym na wykładzie. Czy przy pomocy MPI możliwe jest osiągnięcie maksymalnej 
teoretycznej przepustowości sieci? Jeżeli nie to wymień przyczyny, które Twoim zdaniem to 
uniemożliwiają.

\chapter{Rozwiązanie}

%\noindent\begin{minipage}{\linewidth}
\lstinputlisting[label=jednowatek, caption=Kod programu]{../main.cpp}
%\end{minipage}

\chapter{Wyniki badań}

\begin{table}[t]
\centering
\csvautotabular{../wyniki_czas.csv}
\caption{Czas w zależności od rozmiaru komunikatu.}
\label{tabela_czas}
\end{table}

\begin{figure}[ht!]
    \centering
	\begin{tikzpicture}
		\begin{axis}[xlabel={Rozmiar komunikatu}, ylabel={Czas}, legend style={legend pos=outer north east}]
		\addplot table [x=wielkosc, y=blokujacy, col sep=comma] {../wyniki_czas.csv};
		\addplot table [x=wielkosc, y=nieblokujacy, col sep=comma, ] {../wyniki_czas.csv};
		\addplot table [x=wielkosc, y=ping, col sep=comma, ] {../wyniki_czas.csv};
		\legend{$blokujący$,$nieblokujący$,$ping$}
		\end{axis}
	\end{tikzpicture}
    \caption{Wykres czasu w zależności od rozmiaru komunikatu.}
    \label{wykres_czas}
\end{figure}

\begin{table}[t]
\centering
\csvautotabular{../wyniki_przepustowosc.csv}
\caption{Czas w zależności od rozmiaru komunikatu.}
\label{tabela_przepustowosc}
\end{table}

\begin{figure}[ht!]
    \centering
	\begin{tikzpicture}
		\begin{axis}[xlabel={Rozmiar komunikatu}, ylabel={Przepustowość}, legend style={legend pos=outer north east}]
		\addplot table [x=wielkosc, y=blokujacy, col sep=comma] {../wyniki_przepustowosc.csv};
		\addplot table [x=wielkosc, y=nieblokujacy, col sep=comma, ] {../wyniki_przepustowosc.csv};
		\legend{$blokujacy$,$nieblokujacy$}
		\end{axis}
	\end{tikzpicture}
    \caption{Wykres czasu w zależności od rozmiaru komunikatu.}
    \label{wykres_przepustowosc}
\end{figure}

%\listoffigures

%\bibliography{references}
%\nocite{*}
%\bibliographystyle{unsrt}

\end{document}
